\chapter{Dinámica espacio--temporal}

\section{Extensión a PDEs y configuración numérica}
Se resolvió el sistema de reacción--difusión con elementos finitos (FEniCS). Condiciones iniciales aleatorias se acotan en
\[
0.8 < c < 1.2, \quad 0.2 < s < 0.6, \quad 0 < i < 0.2.
\]
Parámetros de integración: tiempo final \(T = 2\), paso temporal \(\Delta t = 0.005\), malla uniforme con 200 nodos. Las ecuaciones \eqref{eqn:cancer_dynamic_spatial}--\eqref{eqn:inmune_dynamic_spatial} se resolvieron para $\mu\in\{0,1\}$.

\section{Formación de patrones (efecto Allee débil)}
\captionsetup[subfigure]{skip=2pt}
\begin{figure}[H]
    \centering
    \begin{subfigure}[b]{\textwidth}
        \centering
        \includegraphics[width=\textwidth, height=0.28\textheight]{images/fields_mu_0.png}
        \caption{Evolución espacial en \( T = 1.935 \) para \(\mu = 0\).}
    \end{subfigure}
    \begin{subfigure}[b]{\textwidth}
        \centering
        \includegraphics[width=\textwidth, height=0.28\textheight]{images/fields_mu_1.png}
        \caption{Evolución espacial en \( T = 1.935 \) para \(\mu = 1\).}
    \end{subfigure}
    \caption{Distribución espacial de $c$, $s$ e $i$ a \(T=1.935\). Con $\mu=0$ se observan cúmulos caóticos; con $\mu=1$ la estructura variacional suaviza y estabiliza patrones, generando cúmulos más definidos y menos numerosos. Animaciones en \href{https://drive.google.com/drive/folders/11FhjobQToXXDddC7UCPAwGjC53Cm1i0S?usp=sharing}{enlace}.}
    \label{fig:fields_evolution}
\end{figure}
El cáncer domina espacialmente, las células sanas disminuyen y el sistema inmune responde de forma localizada; $\mu$ modula estabilidad y organización.

\section{Espectros de potencia y correlaciones (débil)}
\subsection{Espectros de potencia del campo $c$}
\begin{figure}[H]
    \centering
    \begin{subfigure}[b]{0.45\textwidth}
        \centering
        \includegraphics[width=\linewidth]{images/power_spectrum_c_0.000.png}
        \caption*{\(\mu = 1,\ T = 0\)}
    \end{subfigure}
    \hfill
    \begin{subfigure}[b]{0.45\textwidth}
        \centering
        \includegraphics[width=\linewidth]{images/power_spectrum_c_1.000.png}
        \caption*{\(\mu = 1,\ T = 1\)}
    \end{subfigure}
    \caption{Transición de espectro amplio (desorden) a bajas frecuencias (compactación) para $c$ con $\mu=1$.}
    \label{fig:power_spectrum_cancer_mu1}
\end{figure}

\subsection{Autocorrelaciones y correlaciones cruzadas}
La longitud de correlación $\xi(t)$ sigue $\xi(t)\sim e^{\alpha} t^{0.5}$, con $\alpha$ dependiente de la correlación y de $\mu$. Figuras:
\begin{figure}[H]
    \centering
    \includegraphics[width=0.65\textwidth]{images/correlation_mean_c_c.png}
    \caption{Autocorrelación $c_c$ (células cancerosas). $\mu=1$ incrementa $\xi(t)$, indicando mayor compactación.}
    \label{fig:correlation_c_c}
\end{figure}
\begin{figure}[H]
    \centering
    \includegraphics[width=0.65\textwidth]{images/correlation_mean_s_s.png}
    \caption{Autocorrelación $s_s$ (células sanas). $\mu=1$ produce distribución más homogénea.}
    \label{fig:correlation_s_s}
\end{figure}
\begin{figure}[H]
    \centering
    \includegraphics[width=0.65\textwidth]{images/correlation_mean_i_i.png}
    \caption{Autocorrelación $i_i$ (sistema inmune). Diferencias leves entre $\mu=0$ y $\mu=1$.}
    \label{fig:correlation_i_i}
\end{figure}
\begin{figure}[H]
    \centering
    \includegraphics[width=0.65\textwidth]{images/correlation_mean_c_s.png}
    \caption{Correlación cruzada $c_s$ (cáncer--sanas). Crecimiento débil de $\xi(t)$, coherente con competencia.}
    \label{fig:correlation_c_s}
\end{figure}
\begin{figure}[H]
    \centering
    \includegraphics[width=0.65\textwidth]{images/correlation_mean_c_i.png}
    \caption{Correlación cruzada $c_i$ (cáncer--inmune). $\mu=1$ mejora seguimiento inmune de patrones tumorales.}
    \label{fig:correlation_c_i}
\end{figure}
\begin{figure}[H]
    \centering
    \includegraphics[width=0.65\textwidth]{images/correlation_mean_s_i.png}
    \caption{Correlación cruzada $s_i$ (sanas--inmune). Es la más débil; el sistema inmune se concentra en el tumor.}
    \label{fig:correlation_s_i}
\end{figure}
En conjunto, $\mu=1$ aumenta la coherencia espacial de $c$ y mejora el acoplamiento $c_i$; $s_i$ permanece bajo.

\section{Efecto Allee fuerte: dinámica espacio--temporal}
El efecto fuerte introduce barrera crítica; la dinámica espacial resulta:
\begin{figure}[H]
    \centering
    \begin{subfigure}[b]{1\textwidth}
        \centering
        \includegraphics[width=\textwidth, height=0.3\textheight]{images/fields_strong_mu_0.png}
        \caption{$T = 1.935$, $\mu = 0$.}
    \end{subfigure}
    \vspace{0.2em}
    \begin{subfigure}[b]{1\textwidth}
        \centering
        \includegraphics[width=\textwidth, height=0.3\textheight]{images/fields_strong_mu_1.png}
        \caption{$T = 1.935$, $\mu = 1$.}
    \end{subfigure}
    \caption{Distribución espacial bajo efecto Allee fuerte. $\mu=0$: patrones fragmentados; $\mu=1$: dominios más definidos.}
    \label{fig:fields_strong}
\end{figure}
Con $\mu=0$ hay cúmulos dispersos; con $\mu=1$ se observa segregación más clara y respuesta inmune organizada.

\subsection{Correlaciones (fuerte)}
La evolución de $\xi(t)$ es subdifusiva en todas las correlaciones:
\begin{figure}[H]
    \centering
    \includegraphics[width=0.65\textwidth]{images/correlation_strong_mean_c_c.png}
    \caption{Autocorrelación $c_c$ con efecto Allee fuerte.}
    \label{fig:correlation_strong_c_c}
\end{figure}
\begin{figure}[H]
    \centering
    \includegraphics[width=0.65\textwidth]{images/correlation_strong_mean_s_s.png}
    \caption{Autocorrelación $s_s$ (fuerte).}
    \label{fig:correlation_strong_s_s}
\end{figure}
\begin{figure}[H]
    \centering
    \includegraphics[width=0.65\textwidth]{images/correlation_strong_mean_i_i.png}
    \caption{Autocorrelación $i_i$ (fuerte).}
    \label{fig:correlation_strong_i_i}
\end{figure}
\begin{figure}[H]
    \centering
    \includegraphics[width=0.65\textwidth]{images/correlation_strong_mean_c_s.png}
    \caption{Correlación $c_s$ (fuerte).}
    \label{fig:correlation_strong_c_s}
\end{figure}
\begin{figure}[H]
    \centering
    \includegraphics[width=0.65\textwidth]{images/correlation_strong_mean_c_i.png}
    \caption{Correlación $c_i$ (fuerte).}
    \label{fig:correlation_strong_c_i}
\end{figure}
\begin{figure}[H]
    \centering
    \includegraphics[width=0.65\textwidth]{images/correlation_strong_mean_s_i.png}
    \caption{Correlación $s_i$ (fuerte).}
    \label{fig:correlation_strong_s_i}
\end{figure}
Observaciones: $c_c$ crece lentamente y similar para $\mu=0,1$; $s_s$ es ligeramente mayor con $\mu=0$; $i_i$ es bajo y estable; $c_s$ y $s_i$ muestran acoplamiento débil; $c_i$ permanece débil, reflejando limitada respuesta inmune cuando la barrera Allee domina.

\section{Síntesis}
El efecto Allee débil permite patrones amplios; $\mu=1$ aumenta coherencia y facilita contención parcial. En el efecto fuerte, la barrera crítica limita el tamaño de dominios tumorales y atenúa el impacto de $\mu$, aunque sigue actuando como organizador espacial secundario.




